\documentclass[11pt]{article}
\usepackage[T1]{fontenc}
\usepackage[utf8]{inputenc}
\usepackage{lmodern,microtype}
\usepackage[margin=1in]{geometry}
\usepackage{amsmath,amssymb,amsthm,mathtools,bm}
\usepackage{booktabs,array}
\usepackage{xurl}
\usepackage[hidelinks]{hyperref}
\hypersetup{pdftitle={Equality, rigidity and full-ellipsoid CHSH envelopes at fixed steering-ellipsoid centre},pdfauthor={Dev D. Goyal}}
\newtheorem{theorem}{Theorem}[section]
\newtheorem{lemma}[theorem]{Lemma}
\newtheorem{corollary}[theorem]{Corollary}
\newtheorem{proposition}[theorem]{Proposition}
\theoremstyle{remark}\newtheorem{remark}[theorem]{Remark}
\DeclareMathOperator{\tr}{tr}\DeclareMathOperator{\diag}{diag}
\DeclareMathOperator{\spec}{spec}
\newcommand{\R}{\mathbb R}\newcommand{\C}{\mathbb C}
\newcommand{\id}{\mathbb I}\newcommand{\ket}[1]{|#1\rangle}
\newcommand{\bra}[1]{\langle#1|}\newcommand{\proj}[1]{\ket{#1}\bra{#1}}
\newcommand{\norm}[1]{\left\lVert#1\right\rVert}
\newcommand{\cv}{\bm c}\newcommand{\av}{\bm a}\newcommand{\bv}{\bm b}
\newcommand{\eps}{\varepsilon}
\newcommand{\Dp}{D_{\mathrm{pur}}}
\title{Equality, rigidity and full-ellipsoid CHSH envelopes\protect\\
at fixed steering-ellipsoid centre}
\author{Dev D. Goyal}
\date{3 November 2026}
\begin{document}
\maketitle
\begin{abstract}
Milne, Jennings, Jevtic and Rudolph conjectured that a two-qubit state whose
steering-ellipsoid centre has length $c$ has CHSH value at most
$\max\{2\sqrt{2(1-c)},2\}$ (arXiv:1404.3951, Conjecture~1). The Omega
Institute's public repository recorded a claimed Lean proof of that conjecture
on 14 October 2026, and its proof already contains the refined bound
$\beta^2/4\le 1+(1-r^2)(1-2c)$, where $r$ is the length of Bob's Bloch vector;
an attaining family with this value goes back to Zhang, Araiza Bravo, Lorenz
and Chitambar (2020). We claim no priority for either. Building on them, we prove the
complete equality classification at every fixed pair $(c,r)$: for $0<c<1$ the
extremizers form a single local-unitary class, a filtered maximally obese state
with the filter condition $\norm{\bv}^2=b_z$. We give a quantitative trace-norm
rigidity theorem, including an estimate that preserves the whole centre vector
with exponent $1/4$, and a separable rank-two family showing that this exponent
cannot be improved uniformly in $c$. Beyond the centre alone, we determine the
exact fixed-radius CHSH envelope for every axial steering ellipsoid, with all
equality directions and a resulting exact CHSH-breaking criterion for axial
qubit channels, and for every centred ellipsoid. We also correct a printed
optimizing direction in the 2020 paper. A Lean~4 development checks, for every
admissible $(c,r)$, an actual complex density matrix with the prescribed centre
and Bob length whose CHSH supremum over all unit Pauli settings is exactly
$2\sqrt{1+(1-2c)(1-r^2)}$, two further explicit witnesses, and the scalar
spectral optimization behind the axial envelope; the classification, rigidity
and general filtering steps are conventional proofs.
\end{abstract}

\section{Introduction}

For a two-qubit state, the steering ellipsoid of Alice is the set of Bloch
vectors to which Bob can steer her by local measurements and postselection.
Milne, Jennings, Jevtic and Rudolph~\cite{milne} proved that, for states whose
Bob marginal is maximally mixed, the CHSH value is at most
$2\sqrt{2(1-c)}$, where $c$ is the length of the ellipsoid's centre, and they
conjectured (Conjecture~1 of~\cite{milne}) that every two-qubit state obeys
\begin{equation}\label{eq:milne}
 \beta(\rho)\le\max\{2\sqrt{2(1-c)},2\}.
\end{equation}
Ducuara, Susa and Skrzypczyk~\cite{ducuara} use centre bounds of this kind
in their criterion for locally inaccessible hidden correlations.

\paragraph{Earlier public work and priority (note dated 3 November 2026).}
The Omega Institute's public repository
\texttt{the-omega-institute/trureturing} records a claimed Lean proof of
Conjecture~1 for every two-qubit density matrix, with supremal sharpness at
every centre length; the dossier for its issue \#16558 was committed on
14 October 2026 at 08:40:22 UTC~\cite{dossier}. Its proof combines the inequalities
$M\le1+2\Delta(\rho)$ and $\Delta(\rho)\le(1-r^2)(\tfrac12-c)$ (notation below),
so it proves the refined joint inequality
\begin{equation}\label{eq:H}
 M(\rho):=\frac{\beta(\rho)^2}{4}\le H(c,r):=1+(1-r^2)(1-2c),
\end{equation}
and the dossier states this inequality explicitly as proved on the way.
The same record proves Conjecture~2 of~\cite{milne} (issue \#14206, committed
on 8 October 2026). Independently of both, Appendix~B of Zhang, Araiza Bravo,
Lorenz and Chitambar~\cite{zhang} computes the optimum $1+\alpha^2(2p-1)$ of
CHSH over an amplitude-damping family; with $p=1-c$ and $\alpha^2=1-r^2$ this
is $H(c,r)$. Combining the two, the exact supremum of $\beta$ at fixed $(c,r)$
is $2\sqrt{H(c,r)}$. \emph{We claim no priority for Conjecture~1, for the
joint bound~\eqref{eq:H}, for its attaining value, or for the attaining family.}
Section~\ref{sec:prior} reconstructs the bound only to fix notation and to
extract the slack that the new results use. We have not compiled the Omega
Lean proof and describe it as claimed; its mathematics was read and
reconstructed.

\paragraph{What is new here.}
The following go beyond Conjecture~1 and the Omega record. The Omega dossier
lists bounds involving the full steering ellipsoid as open and argues
attainment only for the centre-only bound.
\begin{enumerate}
\item \emph{Equality classification} (Theorem~\ref{thm:joint}). For $0<c<1$
and $r<1$, $M=H(c,r)$ holds exactly for one local-unitary class: a maximally
obese canonical state filtered by $\sqrt{\id+\bv\cdot\bm\sigma}$ with
$\norm{\bv}^2=b_z$. The necessity step is new; earlier work exhibits
attaining states but does not exclude others. Equality in~\eqref{eq:milne} is
classified for every $c$ (Corollary~\ref{cor:milne-equality}).
\item \emph{Rigidity} (Section~\ref{sec:rigidity}). Near-equality forces
closeness, in trace norm, to an exact extremizer with the same $c$ and $r$;
a version preserves the whole centre vector. Uniformly in $c$, canonical
states lie within $\min\{2,9\Dp^{1/4}\}$ of a maximally obese state with the
same Alice marginal, and the exponent $1/4$ is optimal.
\item \emph{Full-ellipsoid envelopes} (Section~\ref{sec:shapes}). Exact
fixed-radius CHSH maxima, with all optimizing filter directions, for every
translated sphere, every axial ellipsoid (centre on the symmetry axis, oblate
or prolate) and every centred ellipsoid, and the closed optimization over the
radius. For axial qubit channels $x\mapsto ce_z+\diag(s,s,t)x$ this gives an
exact criterion for CHSH breaking.
\item \emph{A printed optimizer} (Section~\ref{sec:zhang}). The optimizing
direction displayed in Appendix~B of~\cite{zhang} is incorrect; an exact
example shows it does not attain the stated maximum. The maximal value and
threshold of~\cite{zhang} are unaffected.
\end{enumerate}
The literature comparison is bounded; it does not establish worldwide
novelty.

\paragraph{Verification.}
A standalone Lean~4.19.0 project (Section~\ref{sec:formal}) proves, without
assuming the Horodecki formula, that an explicit complex $4\times4$ density
matrix with centre $(0,0,c)$ and Bob length $r$ has CHSH supremum exactly
$2\sqrt{H(c,r)}$ for every admissible $(c,r)$, and it checks the axial witness
of Section~\ref{sec:axial-witness}, the sharpness family of
Corollary~\ref{cor:uniform-rigidity} and the scalar spectral optimization of
Theorem~\ref{thm:axial}. The general filtering correspondence, the equality
classification and the rigidity theorems are conventional proofs, supported
by exact symbolic identities and seeded numerical falsification searches.

\section{Definitions and conventions}\label{sec:defs}

A two-qubit density matrix has the Bloch representation
\begin{equation}\label{eq:bloch}
 \rho=\frac14\Bigl(\id\otimes\id+\av\cdot\bm\sigma\otimes\id
 +\id\otimes\bv\cdot\bm\sigma+\sum_{i,j=1}^{3}T_{ij}\sigma_i\otimes\sigma_j\Bigr),
\end{equation}
with Alice the first factor. Put $r=\norm{\bv}$. Following~\cite{milne}, four
unit vectors define
\[
 B=(\alpha\cdot\bm\sigma)\otimes((\beta+\beta')\cdot\bm\sigma)
 +(\alpha'\cdot\bm\sigma)\otimes((\beta-\beta')\cdot\bm\sigma),
 \qquad \beta(\rho)=\max_B|\tr(\rho B)|.
\]
Constant observables $\pm\id$ are not allowed, so $\beta$ is not floored at
two. By the Horodecki theorem~\cite{horodecki},
\begin{equation}\label{eq:horodecki}
 M(\rho)=\frac{\beta(\rho)^2}{4}
 =\lambda_1(TT^{\mathsf T})+\lambda_2(TT^{\mathsf T})
 =\max_{u\perp v,\ \norm u=\norm v=1}\bigl(\norm{T^{\mathsf T}u}^2+\norm{T^{\mathsf T}v}^2\bigr).
\end{equation}
For $r<1$, Alice's steering ellipsoid has centre and canonical form
\begin{equation}\label{eq:centre}
 \cv=\frac{\av-T\bv}{1-r^2},\quad c=\norm{\cv},\qquad
 \sigma=(\id\otimes(2\rho_B)^{-1/2})\rho(\id\otimes(2\rho_B)^{-1/2}),
 \quad\sigma_B=\id/2.
\end{equation}
The canonical state $\sigma$ has Alice Bloch vector $\cv$ and correlation
matrix $C$; the shape matrix is $Q=CC^{\mathsf T}$, whose eigenvalues are the
squared semiaxes. For $r=1$ positivity forces $\rho=\rho_A\otimes\rho_B$, the
ellipsoid is the point $\{\av\}$ and we set $\cv=\av$, the convention also
used in~\cite{dossier}.

\paragraph{The centre formula.}
For a Bob effect proportional to $\id+e\cdot\bm\sigma$, $\norm e\le1$, Alice's
conditional Bloch vector is $(\av+Te)/(1+\bv\cdot e)$. An invertible positive
Bob filter $F$ maps effects bijectively, up to positive rescaling, by
$E\mapsto FEF$, so whitening preserves the steering set. Put $s=\sqrt{1-r^2}$,
$W=s\id+\bv\bv^{\mathsf T}/(1+s)$ and $F=\sqrt{\id+\bv\cdot\bm\sigma}$. Pauli
multiplication gives $F(x\cdot\bm\sigma)F=(\bv\cdot x)\id+(Wx)\cdot\bm\sigma$,
$W\bv=\bv$ and $W^2=s^2\id+\bv\bv^{\mathsf T}$. If the canonical state has Alice
vector $a_0$ and correlation matrix $C$, undoing the whitening gives
$\av=a_0+C\bv$, $T=a_0\bv^{\mathsf T}+CW$ and $\av-T\bv=(1-r^2)a_0$. The
canonical steering set is $\{a_0+Ce:\norm e\le1\}$, with centre $a_0$ and shape
matrix $CC^{\mathsf T}$; hence $a_0=\cv$. The same identities give the Gram
formula
\begin{equation}\label{eq:gram}
 \av=\cv+C\bv,\qquad
 TT^{\mathsf T}=\av\av^{\mathsf T}+(1-r^2)(Q-\cv\cv^{\mathsf T}),
\end{equation}
valid for a state $\rho=(\id\otimes F_{\bv})\sigma(\id\otimes F_{\bv})$ obtained
from a canonical state by the filter $F_{\bv}=\sqrt{\id+\bv\cdot\bm\sigma}$.
Every state with $r<1$ has this form up to a local unitary on Bob, since
$\rho=(\id\otimes\sqrt{2\rho_B})\sigma(\id\otimes\sqrt{2\rho_B})$.

\section{The joint bound (prior work)}\label{sec:prior}

This section reconstructs~\eqref{eq:H} along the route of the Omega
proof~\cite{dossier}; nothing in it is claimed as new. Define
$\Delta(X)=\tr(X^2)-\tr\bigl((\tr_BX)^2\bigr)$.

\begin{lemma}[CHSH--purity bound]\label{lem:purity}
For every two-qubit density matrix, $M(\rho)\le1+2\Delta(\rho)$.
\end{lemma}
\begin{proof}
Let $u,v,n$ be an orthonormal frame, $K_n=\sigma_y\otimes((n\cdot\bm\sigma)\sigma_y)$
and $\rho_n=K_n\rho^{\mathsf T}K_n^\dagger\succeq0$ (full transpose). The map
$\rho\mapsto\rho_n$ sends $(\av,\bv,T)$ to
$(-\av,\bv-2(\bv\cdot n)n,T(2nn^{\mathsf T}-\id))$, so Pauli orthogonality gives
$4\tr(\rho\rho_n)=1-\norm{\av}^2+\norm{\bv}^2-2(\bv\cdot n)^2+2\norm{Tn}^2-\norm{T}_F^2$.
With $4\Delta=\norm{T}_F^2+\norm{\bv}^2-\norm{\av}^2-1$,
\begin{equation}\label{eq:flip}
 \norm{Tu}^2+\norm{Tv}^2+(\bv\cdot n)^2+2\tr(\rho\rho_n)=1+2\Delta(\rho).
\end{equation}
The last two terms on the left are nonnegative; maximize over the plane.
\end{proof}

\begin{lemma}[Filter covariance]\label{lem:covariance}
For every $4\times4$ matrix $X$ and every complex $2\times2$ matrix $F$,
$\Delta((\id\otimes F)X(\id\otimes F)^\dagger)=|\det F|^2\Delta(X)$.
\end{lemma}
\begin{proof}
For $2\times2$ matrices, $\tr(AB)-\tr A\tr B=\det A+\det B-\det(A+B)$, and
simultaneous congruence by $F$ multiplies the right side by $|\det F|^2$.
Apply this to the Bob blocks of $X$ and sum over Alice indices.
\end{proof}

\begin{lemma}[Canonical caps]\label{lem:cap}
If $\sigma_B=\id/2$ and Alice's Bloch length is $c$, then
$\lambda_{\max}(\sigma)\le L:=1-c/2$, $\tr\sigma^2\le L^2+(1-L)^2=1-c+c^2/2$
and $\Delta(\sigma)\le1/2-c$.
\end{lemma}
\begin{proof}
Let $\lambda$ be the largest eigenvalue with unit eigenvector $\psi$ and
$\tau=\sigma-\lambda\proj\psi\succeq0$. Bob's marginal gives
$\lambda\psi_B\preceq\id/2$; the marginals of a pure state have equal spectra,
so $\lambda\psi_A\preceq\id/2$. For a top eigenvector $e_+$ of $\sigma_A$,
$1-\lambda=\tr\tau_A\ge\bra{e_+}\tau_A\ket{e_+}=\tfrac{1+c}2-\lambda\bra{e_+}\psi_A\ket{e_+}\ge c/2$.
If $\lambda\le1/2$ then $\tr\sigma^2\le\lambda\le1/2\le L^2+(1-L)^2$; otherwise
$\tr\sigma^2\le\lambda^2+(1-\lambda)^2$, increasing for $\lambda\ge1/2$.
Subtract Alice's purity $(1+c^2)/2$.
\end{proof}

For $r<1$ write $\rho=(\id\otimes F)\sigma(\id\otimes F)$ with
$F=\sqrt{2\rho_B}$ and $|\det F|^2=1-r^2$; the three lemmas give~\eqref{eq:H}.
If $r=1$ the state is a product and $M=c^2\le1=H$. Maximizing $H$ over $r$
gives~\eqref{eq:milne}.

\section{Canonical equality and a shape-sensitive deficit}

Define the canonical maximally obese state
\begin{equation}\label{eq:obese}
 \omega_c=\proj{w_c}+\frac c2\proj{01},\qquad
 \ket{w_c}=\frac{\ket{00}+\sqrt{1-c}\ket{11}}{\sqrt2},
\end{equation}
with $\cv=ce_z$, $C=\diag(\sqrt{1-c},-\sqrt{1-c},1-c)$ and
$\spec(\omega_c)=\{1-c/2,c/2,0,0\}$. It is locally unitarily equivalent to
the maximally obese state of~\cite{milne}.

\begin{theorem}[Canonical equality]\label{thm:canonical-equality}
Equality in the spectral cap or in the purity cap of Lemma~\ref{lem:cap} holds
if and only if $\sigma$ is a local-unitary image of $\omega_c$.
\end{theorem}
\begin{proof}
Let $0<c\le1$ and $\lambda=L$. Every inequality in the proof of the spectral
cap is an equality, so $\tau_A=\frac c2\proj{e_+}$ and
$\lambda\psi_A=\frac12\proj{e_+}+\frac{1-c}2\proj{e_-}$. A positive operator
with a rank-one marginal is a product on that factor, so
$\tau=\proj{e_+}\otimes\tau_B$. In Schmidt bases
$\sqrt\lambda\ket\psi=\ket{e_+,f_0}/\sqrt2+\sqrt{(1-c)/2}\,\ket{e_-,f_1}$, and
Bob's marginal forces $\tau_B=\frac c2\proj{f_1}$: this is~\eqref{eq:obese}
in local bases. At $c=0$, $\lambda=1$ gives a pure state with maximally mixed
Bob marginal, a maximally entangled state. The converse is immediate. For
$c<1$ the purity cap exceeds $1/2$, so equality in it forces $\lambda>1/2$ and,
by strict monotonicity of $x^2+(1-x)^2$, $\lambda=L$. At $c=1$ every canonical
state is $\proj{e_+}\otimes\id/2$.
\end{proof}

Define the nonnegative shape deficit
\begin{equation}\label{eq:D}
 \Dp=1-c+\frac{c^2}2-\tr\sigma^2=\frac{3-4c+c^2-\tr Q}{4},
\end{equation}
using $4\tr\sigma^2=1+c^2+\tr Q$. Theorem~\ref{thm:canonical-equality} says
$\Dp=0$ exactly on the obese orbit. Keeping this slack in the chain of
Section~\ref{sec:prior} gives
\begin{equation}\label{eq:shape-gap}
 M(\rho)\le H(c,r)-2(1-r^2)\Dp=1+\frac{1-r^2}{2}\bigl(\tr Q-c^2-1\bigr).
\end{equation}
This is an immediate refinement of the prior chain, not a sharp bound for a
general ellipsoid, since orientation is discarded; it penalizes every
canonical state that is not maximally obese.

\section{Equality in the joint bound}\label{sec:joint}

\begin{theorem}[Equality classification]\label{thm:joint}
For $0\le c\le1$ and $0\le r<1$,
\begin{equation}\label{eq:joint-max}
 \max_{\rho:\ \norm{\cv}=c,\ \norm{\bv}=r}\beta(\rho)=2\sqrt{1+(1-r^2)(1-2c)}.
\end{equation}
For $0<c<1$ the attaining states are, up to local unitaries, exactly
\begin{equation}\label{eq:rho-b}
 \rho_{c,\bv}=(\id\otimes F_{\bv})\omega_c(\id\otimes F_{\bv}),
 \qquad F_{\bv}=\sqrt{\id+\bv\cdot\bm\sigma},\qquad \norm{\bv}^2=b_z<1,
\end{equation}
for instance $\bv=(r\sqrt{1-r^2},0,r^2)$. All transverse azimuths are
equivalent under local unitaries fixing $\omega_c$, so each interior pair
$(c,r)$ has one local-unitary class of extremizers. At $c=0$ the extremizers
are the pure states with Bob length $r$; at $c=1$ every state is a product
with pure Alice marginal and attains $2r$.
\end{theorem}
The value~\eqref{eq:joint-max} is prior work (Section~1); the classification
is the new content.
\begin{proof}
The upper bound is~\eqref{eq:H}. If equality holds with $r<1$, every
inequality in $M\le1+2\Delta(\rho)=1+2(1-r^2)\Delta(\sigma)\le H$ is an
equality, so $\Dp=0$, and Theorem~\ref{thm:canonical-equality} reduces every
extremizer to a filtered $\omega_c$ (the unitary freedom in $F$ is absorbed into
local unitaries). It remains to find the filters that attain the cap. Let
$R=\norm{\bv}^2$, $p=1-c$, and rotate so that $\bv=(\sqrt u,0,z)$, $u\ge0$,
$R=u+z^2<1$. By~\eqref{eq:gram}, with
$q=p(1-R)$, $k=(1-2c)(1-R)$ and $h=1+k-q=p+cR$, the $y$ direction is an
eigenvector of $TT^{\mathsf T}$ with eigenvalue $q$, and the $xz$ block is
\begin{equation}\label{eq:block}
 B=\begin{pmatrix}
 p(1-z^2)&\sqrt{pu}(c+pz)\\
 \sqrt{pu}(c+pz)&k+(c+pz)^2
 \end{pmatrix}.
\end{equation}
Its eigenvalues $\lambda_-\le\lambda_+$ satisfy $\lambda_-\le q\le\lambda_+$,
because $\det(q\id-B)=-cp\,u(1-R)\le0$; hence $M=q+\lambda_+$. Moreover
\begin{equation}\label{eq:factor}
 (h\id-B)_{11}=cu+z^2,\qquad (h\id-B)_{22}=p\{u+c(z-1)^2\},\qquad
 \det(h\id-B)=cp(R-z)^2,
\end{equation}
so $h\id-B\succeq0$ and $M\le q+h=H$. For $0<c<1$ equality holds exactly when
$\det(h\id-B)=0$, that is $R=z$. This gives both directions of the
classification; the displayed $\bv$ has $R=r^2=z$. The filter is invertible
and positive and yields trace one because $(\omega_c)_B=\id/2$; scaled by
$1/\sqrt{1+r}$ it is a valid Kraus operator. At $c=0$, $\Dp=0$ forces a
maximally entangled canonical state, and invertible filtering gives exactly the
entangled pure states with $r<1$, with $M=2-r^2$. At $c=1$ a pure Alice
marginal forces a product and $M=r^2$.
\end{proof}
For $0<c<1$ and $r<1$ the classification reads, without coordinates,
\[
 M=H\iff \Dp=0\ \text{and}\ \cv\cdot(\av-\cv)=c(1-c)r^2,
\]
since $\av-\cv=C\bv$ and in the obese frame the left side is $c(1-c)b_z$.

\begin{proposition}[Explicit extremizer]\label{prop:dyadic}
For $0\le c\le1$ and $0\le r<1$ put $p=1-c$, $s=\sqrt{1-r^2}$, $a=\sqrt p\,s$,
\[
 v_0=\tfrac12(1+r,-(1-r),a,a)^{\mathsf T},\qquad
 v_1=\tfrac{\sqrt c\,s}{2}(1,1,0,0)^{\mathsf T},\qquad
 \rho^\sharp_{c,r}=v_0v_0^{\mathsf T}+v_1v_1^{\mathsf T}.
\]
Then $\rho^\sharp_{c,r}$ is a density matrix with $\bv=(0,0,r)$,
$\av=(ra,0,c+pr^2)$, $\cv=(0,0,c)$,
\[
 T=\begin{pmatrix}ra&0&a\\0&-a&0\\-p(1-r^2)&0&r\end{pmatrix},\qquad
 \det(x\id-TT^{\mathsf T})=(x-h)(x-q)(x-qh),
\]
with $h=p+cr^2\ge q=p(1-r^2)\ge qh$; consequently $M=h+q=H(c,r)$.
\end{proposition}
\begin{proof}
Positivity is manifest and $\norm{v_0}^2+\norm{v_1}^2=1$. Partial traces give
$\rho_B=\diag((1+r)/2,(1-r)/2)$ and the stated $\av$; Pauli traces give $T$;
$\av-T\bv=(0,0,c(1-r^2))$ gives the centre. The characteristic polynomial and
ordering follow from $a^2=q$, $h-q=r^2$ and $q-qh=qc(1-r^2)$. Whitening gives
$C=\diag(\sqrt p,-\sqrt p,p)R_y$ with a rotation $R_y$, so $Q=\diag(p,p,p^2)$.
\end{proof}
This is a coordinate realization of the attaining family discussed above,
not a new family. Its Bell supremum is the formal endpoint of
Section~\ref{sec:formal}.

\begin{remark}[Pure Bob marginal]
At $r=1$ the constrained maximum is $2c$, not the limit $2$
of~\eqref{eq:joint-max}: the state is a product and its centre is $\av$. The
centre can jump when an invertible filter becomes rank one.
\end{remark}

\begin{corollary}[Equality in Conjecture~1]\label{cor:milne-equality}
Up to local unitaries, equality in~\eqref{eq:milne} holds exactly as follows:
for $0\le c<1/2$, only for $\omega_c$; for $c=1/2$, exactly for the states
$\rho_{1/2,\bv}$ with $R=b_z<1$, whose filter vectors form the sphere
$b_x^2+b_y^2+(b_z-1/2)^2=1/4$ without its north pole; for $1/2<c<1$, never
(the supremum two is not attained); for $c=1$, exactly for pure product states.
\end{corollary}
\begin{proof}
For $c<1/2$, $H=2(1-c)-r^2(1-2c)$ forces $r=0$. For $c=1/2$, $H=1$ for every
$r<1$, and Theorem~\ref{thm:joint} applies. For $1/2<c<1$, $H<1$ when $r<1$,
and a pure Bob marginal gives $M=c^2<1$. For $c=1$, $M=r^2$ requires $r=1$.
\end{proof}
The Omega dossier argues, without formalizing it, that~\eqref{eq:milne} is
attained exactly for $c\in[0,1/2]\cup\{1\}$; the corollary adds the
classification of the attaining states. At $c=1/2$ equality does not force a
maximally mixed Bob marginal: the matrix
\[
 \rho_*=\frac1{32}\begin{pmatrix}
 21&-3&3\sqrt6&3\sqrt6\\
 -3&5&-\sqrt6&-\sqrt6\\
 3\sqrt6&-\sqrt6&3&3\\
 3\sqrt6&-\sqrt6&3&3
 \end{pmatrix}
 =v_0v_0^{\mathsf T}+v_1v_1^{\mathsf T},\quad
 v_0=\bigl(\tfrac34,-\tfrac14,\tfrac{\sqrt6}8,\tfrac{\sqrt6}8\bigr)^{\mathsf T},\
 v_1=\bigl(\tfrac{\sqrt6}8,\tfrac{\sqrt6}8,0,0\bigr)^{\mathsf T},
\]
has $\rho_B=\diag(3/4,1/4)$, $\cv=(0,0,1/2)$,
$\spec(TT^{\mathsf T})=\{5/8,3/8,15/64\}$ and hence $\beta(\rho_*)=2$.

\section{Quantitative rigidity}\label{sec:rigidity}

The norm $\norm{\cdot}_1$ is the full trace norm (twice the trace distance).

\begin{lemma}[Spectral deficit]\label{lem:slack}
Let $0\le c<1$, $p=1-c$, $L=1-c/2$ and $\delta=L-\lambda_{\max}(\sigma)$ for a
canonical $\sigma$. Then $0\le\delta\le\Dp/p$.
\end{lemma}
\begin{proof}
With $m=\lambda_{\max}$ we have $\tr\sigma^2\le c/2+pm$: if $m\le1/2$ use
$\tr\sigma^2\le m$; if $m\ge1/2$ use $\tr\sigma^2\le m^2+(1-m)^2$ and
$c/2+pm-m^2-(1-m)^2=(L-m)(2m-1)\ge0$. Subtract from $L^2+(1-L)^2$ and use
$L^2+(1-L)^2-pL=c/2$.
\end{proof}

\begin{theorem}[Trace-norm rigidity]\label{thm:stability}
If $0\le c<1$ and $\Dp\le(1-c)^2/4$, some local-unitary image $\omega$ of
$\omega_c$ satisfies
\begin{equation}\label{eq:stability}
 \norm{\sigma-\omega}_1
 \le2\sqrt{\frac{\Dp}{1-c}}+\frac{2\sqrt{\Dp}}{1-c}+\frac{\Dp}{1-c}
 \le\frac{9\sqrt{\Dp}}{2(1-c)}.
\end{equation}
The target has the same centre length; the constants are not claimed sharp.
\end{theorem}
\begin{proof}
Let $\lambda=L-\delta$ be the largest eigenvalue, $\sigma=\proj v+\tau$ with
$v=\sqrt\lambda\psi$, $\tau\succeq0$, $\tau v=0$, $d=c/2$ and
$t=\tr\tau=d+\delta$. In Schmidt bases
$v=\sqrt{s_+}\ket{e_0f_0}+\sqrt{s_-}\ket{e_1f_1}$ with $s_+=1/2-\kappa$ and
$s_-=p/2-\delta+\kappa$. Bob's marginal gives $\kappa\ge0$. The largest
eigenvalue of $\sigma_A$ is at most $s_++t$, so
$(1+c)/2\le1/2-\kappa+d+\delta$ and $\kappa\le\delta$. Lemma~\ref{lem:slack}
and the hypothesis give $\delta\le\Dp/p\le p/4$, hence $s_-\ge p/4>0$.

Let $\phi=e_0\otimes f_1$. Bob's marginal gives
$\tau_B=\kappa\proj{f_0}+(t-\kappa)\proj{f_1}$. With
$A=\bra{e_0f_0}\tau\ket{e_0f_0}$ and $B=\bra{e_1f_1}\tau\ket{e_1f_1}$, the
relation $\tau v=0$ gives $s_+A=s_-B$. Since $A\le\kappa$, the weight of $\tau$
outside $\phi$ is $\ell=t-\bra\phi\tau\ket\phi=\kappa+B\le\kappa\lambda/s_-$.
For a spectral decomposition $\tau=\sum_j\mu_j\proj{\chi_j}$, the identity
$\norm{\proj{\chi_j}-\proj\phi}_1=2\sqrt{1-|\langle\phi,\chi_j\rangle|^2}$, the
triangle inequality and Cauchy--Schwarz give
$\norm{\tau-t\proj\phi}_1\le2\sqrt{t\ell}\le2\sqrt{t\lambda\kappa/s_-}\le2\sqrt{\delta/p}$,
using $t+\lambda=1$; replacing $t$ by $d$ costs $\delta$. Let
$w=\ket{e_0f_0}/\sqrt2+\sqrt{p/2}\ket{e_1f_1}$ and
$\omega=\proj w+d\proj\phi$. From $(\sqrt x-\sqrt y)^2\le|x-y|$,
$\norm{v-w}^2\le\kappa+(\delta-\kappa)=\delta$, so
$\norm{\proj v-\proj w}_1\le(\norm v+\norm w)\norm{v-w}\le2\sqrt\delta$.
Altogether $\norm{\sigma-\omega}_1\le2\sqrt\delta+2\sqrt{\delta/p}+\delta$;
then use $\delta\le\Dp/p$, $p\le1$ and $\sqrt{\Dp}\le p/2$.
\end{proof}

\begin{corollary}[Near saturation]\label{cor:near}
Let $\eps=H(c,r)-M(\rho)$ with $c<1$, $r<1$ and $\eps\le(1-r^2)(1-c)^2/2$.
Then some local-unitary image $\omega$ of $\omega_c$ satisfies
$\norm{\sigma-\omega}_1\le\frac{9}{2(1-c)}\sqrt{\eps/(2(1-r^2))}$.
\end{corollary}
This follows from $\Dp\le\eps/[2(1-r^2)]$ in~\eqref{eq:shape-gap}. Undoing the
whitening multiplies the bound by at most $1+r=\norm{F}_{\mathrm{op}}^2$; the
approximant is a filtered obese state, and exact extremality also needs the
orientation condition, which the next results control.

\begin{proposition}[Orientation deficit]\label{prop:orientation}
For a filtered obese state with $0<c<1$ and $R<1$, in the notation
of~\eqref{eq:block},
$(H-M)(h-\lambda_-)=c(1-c)(R-z)^2$, and hence $H-M\ge c(1-c)(R-z)^2$.
\end{proposition}
\begin{proof}
The identity is~\eqref{eq:factor} with $M=q+\lambda_+$ and $H=q+h$. Since
$0\le\lambda_-\le h\le1$, and $h-\lambda_->0$ (if $R>0$ use $h-q=R$ and
$\lambda_-\le q$; if $R=0$ it equals $c(1-c)$), divide.
\end{proof}

\begin{lemma}[Fixed-radius filter transport]\label{lem:filter-metric}
Let $\omega$ be any density matrix with $\omega_B=\id/2$, let
$\rho_{\bv}=(\id\otimes F_{\bv})\omega(\id\otimes F_{\bv})$, and let
$\norm{\bv}=\norm{\bm d}=r<1$. With $\gamma_r=(\sqrt{1+r}+\sqrt{1-r})^{-1}$ and
$x=\norm{\bv-\bm d}$,
\[
 \norm{\rho_{\bv}-\rho_{\bm d}}_1\le2\gamma_rx\sqrt{1-\gamma_r^2x^2/4}\le2\gamma_rx,
\]
with equality in the first bound when $\omega$ is maximally entangled.
\end{lemma}
\begin{proof}
At fixed radius $F_{\bv}=\alpha_r\id+\gamma_r\bv\cdot\bm\sigma$ with
$\alpha_r=(\sqrt{1+r}+\sqrt{1-r})/2$ and $\alpha_r^2+\gamma_r^2r^2=1$. Apply both
filters to a purification of $\omega$; both vectors have norm one because
$\omega_B=\id/2$, and their overlap is
$\frac12\tr(F_{\bv}F_{\bm d})=1-\gamma_r^2x^2/2$. The trace-norm distance of two
pure states is twice the square root of one minus the squared overlap, and the
partial trace contracts the trace norm. For maximally entangled $\omega$ no
environment is traced out.
\end{proof}

\begin{theorem}[Distance to an exact extremizer]\label{thm:physical-rigidity}
Let $0<c<1$, $0\le r<1$, $p=1-c$, $\eps=H(c,r)-M(\rho)$ and
$0\le\eps\le(1-r^2)p^2/2$. Put $d=\frac9{2p}\sqrt{\eps/(2(1-r^2))}$,
$t_0=(1+r)d$ and $J_r=2/(\sqrt{1+r}+\sqrt{1-r})$. Some density matrix $\tau$
with the same $c$ and $r$ and $M(\tau)=H(c,r)$ satisfies
\begin{equation}\label{eq:physical-rigidity}
 \norm{\rho-\tau}_1\le t_0+J_r\sqrt{\frac2{1-r}}\sqrt{\frac{\eps+4t_0}{c(1-c)}}.
\end{equation}
For fixed interior $(c,r)$ this is $O(\eps^{1/4})$; when $r=0$ the bound
$\norm{\rho-\tau}_1\le d$ holds.
\end{theorem}
\begin{proof}
Corollary~\ref{cor:near} gives an obese $\omega$ with
$\norm{\sigma-\omega}_1\le d$, so $\rho'=F_{\bv}\omega F_{\bv}$ has
$\norm{\rho-\rho'}_1\le t_0$. Every CHSH operator has norm at most
$2\sqrt2$ (its square is $4\id-[A,A']\otimes[B,B']$), so
$|M(\rho)-M(\rho')|\le4\norm{\rho-\rho'}_1$ and $H-M(\rho')\le\eps+4t_0$. In the
frame where $\omega=\omega_c$, Proposition~\ref{prop:orientation} gives
$c(1-c)(r^2-b_z)^2\le\eps+4t_0$. Let $\bv^*$ have the radius and azimuth of
$\bv$ and $b^*_z=r^2$. With $y=\sqrt{r^2-b_z^2}$ and $w=r\sqrt{1-r^2}$,
$\norm{\bv-\bv^*}^2=(b_z-r^2)^2+(y-w)^2\le\frac2{1-r}(b_z-r^2)^2$, using
$|b_z+r^2|\le r(1+r)$ and $y+w\ge r\sqrt{1-r^2}$. The state
$\tau=F_{\bv^*}\omega_cF_{\bv^*}$ is an exact extremizer by
Theorem~\ref{thm:joint}; apply Lemma~\ref{lem:filter-metric} with
$2\gamma_r=J_r$.
\end{proof}

\begin{theorem}[Rigidity with the whole centre vector fixed]\label{thm:directional}
Let $\sigma_B=\id/2$, let Alice's Bloch vector be $cn$ with $0<c<1$,
$\norm n=1$, and let $\delta=L-\lambda_{\max}(\sigma)\le c/8$. There is an
obese state $\omega$ with $\omega_A=\sigma_A$, $\omega_B=\id/2$ and
$\norm{\sigma-\omega}_1\le\min\{2,9\sqrt{\delta/c}\}$. In particular, if
$\Dp\le c(1-c)/8$, the bound $9\sqrt{\Dp/[c(1-c)]}$ holds with the same Alice
centre vector.
\end{theorem}
\begin{proof}
Set $q=c/2$, $p=1-c$, $\lambda=L-\delta$ and $t_1=1-\lambda=q+\delta$. Write
$\sigma=\lambda\proj\psi+t_1\eta$ with $\eta$ a density matrix, let
$a_\psi,b_\psi,a_\eta,b_\eta$ be the marginal Bloch vectors, $s=\norm{a_\psi}=\norm{b_\psi}$
and $u=\norm{b_\eta}$. Bob-flatness gives $\lambda s=t_1u$ with opposite Bob
directions, and Alice's marginal gives the nonnegative decomposition
\begin{equation}\label{eq:directional-deficit}
 2\delta=t_1(1-u)+\lambda(s-n\cdot a_\psi)+t_1(1-n\cdot a_\eta).
\end{equation}
In particular $t_1u\ge q-\delta>0$. Let $e_0$ point along $n$, $f_1$ along
$b_\eta/u$, and $\Phi=\proj{e_0f_1}$. The two local failure projectors commute,
so \eqref{eq:directional-deficit} gives $\bra{e_0f_1}\eta\ket{e_0f_1}\ge1-\delta/t_1$
and $\norm{t_1\eta-t_1\Phi}_1\le2\sqrt{t_1\delta}$. If $\theta$ is the angle
between $a_\psi/s$ and $n$, the middle term gives
$\sin^2(\theta/2)\le\delta/(\lambda s)\le\delta/(q-\delta)$. Rotating Alice's
Schmidt axis to $n$ about an axis perpendicular to $a_\psi$ changes the vector
by overlap $\cos(\theta/2)$, at trace-norm cost at most
$2\lambda\sqrt{\delta/(q-\delta)}$. With $\kappa=t_1(1-u)/2\in[0,\delta]$ the
rotated weighted vector is
$\sqrt{1/2-\kappa}\ket{e_0f_0}+\sqrt{p/2-\delta+\kappa}\ket{e_1f_1}$; compare it
with $w=\ket{e_0f_0}/\sqrt2+\sqrt{p/2}\ket{e_1f_1}$ and set
$\omega=ww^\dagger+q\Phi$, which is obese with the required marginals. The
vector error is at most $\sqrt\delta$ and the rank-one error at most
$(\sqrt\lambda+\sqrt L)\sqrt\delta$; changing the weight $t_1$ to $q$ costs
$\delta$. Hence
$\norm{\sigma-\omega}_1\le2\sqrt{t_1\delta}+2\lambda\sqrt{\delta/(q-\delta)}
+(\sqrt\lambda+\sqrt L)\sqrt\delta+\delta\le9\sqrt{\delta/c}$ for
$\delta\le c/8$. Lemma~\ref{lem:slack} gives the final statement.
\end{proof}
Consequently Theorem~\ref{thm:physical-rigidity} holds for an exact extremizer
with the same centre \emph{vector} and Bob length under
$\eps\le(1-r^2)c(1-c)/4$, with $d$ replaced by
$9\sqrt{\eps/(2(1-r^2)c(1-c))}$; Bob filtering preserves the steering
ellipsoid, so the orientation correction keeps the centre vector.

\begin{lemma}[Small-centre direction correction]\label{lem:small-centre}
For $0\le c\le1/2$, any local-unitary image of $\omega_c$ can be moved within
its orbit to any prescribed Alice centre direction at trace-norm cost at most
$3c$.
\end{lemma}
\begin{proof}
With $\Phi$ the projector onto $(\ket{00}+\ket{11})/\sqrt2$,
$\norm{\omega_c-\Phi}_1\le\sqrt2(1-\sqrt{1-c})+c/2$ in every local frame. An
Alice rotation paired with $U\otimes\overline U$ fixes $\Phi$, so the cost is at
most $2\sqrt2(1-\sqrt{1-c})+c=\bigl(\tfrac{2\sqrt2}{1+\sqrt{1-c}}+1\bigr)c\le3c$.
\end{proof}

\begin{corollary}[Uniform rigidity with optimal exponent]\label{cor:uniform-rigidity}
For every canonical state, including $c=0,1$, there is an obese state $\omega$
with the same Alice Bloch vector and
\begin{equation}\label{eq:uniform-quarter}
 \norm{\sigma-\omega}_1\le\min\{2,9\Dp^{1/4}\}.
\end{equation}
No bound $K\Dp^{\alpha}$ with $\alpha>1/4$ and a constant $K$ independent of $c$
holds for all canonical states. The constant nine is not claimed optimal.
\end{corollary}
\begin{proof}
\emph{Boundary estimate.} In Alice's eigenbasis write
$\sigma=\left(\begin{smallmatrix}A&B\\B^\dagger&E\end{smallmatrix}\right)$ with
$A+E=\id/2$ and $\tr E=p/2$, $p=1-c$. Positivity gives $B=A^{1/2}KE^{1/2}$ with
a contraction $K$, hence $\norm B_1\le\sqrt{\tr A\tr E}$, and with
$P=\proj0\otimes\id/2$ we get $\norm{\sigma-P}_1\le p+\sqrt{2p-p^2}$. The
aligned obese state obeys the same bound, so the distance to an obese state
with the same Alice vector is at most $2p+2\sqrt{2p-p^2}\le(2+2\sqrt2)\sqrt p$.

\emph{Cases.} Assume $0<\Dp\le1/256$ and put $d=\Dp^{1/4}\le1/4$ (for
$\Dp>1/256$ the bound exceeds two). If $c<1/2$ then $p>1/2$ and
Theorem~\ref{thm:stability} gives distance at most $9\Dp^{1/2}=9d^2$ to some
obese state; if moreover $c<d$, Lemma~\ref{lem:small-centre} fixes the
direction at total cost $9d^2+3d\le9d$, while if $d\le c<1/2$ then
$\Dp=d^4\le d/16\le cp/8$ and Theorem~\ref{thm:directional} gives at most
$9\sqrt2d^{3/2}\le9d$. If $c\ge1/2$ and $p<2d^2$, the boundary estimate gives
at most $(4+2\sqrt2)d<9d$. If $c\ge1/2$ and $p\ge2d^2$, then
$\Dp\le d^2/8\le cp/8$ and $cp\ge d^2$, and Theorem~\ref{thm:directional} gives
at most $9d$.

\emph{Optimality.} For $0<c<1$ let
\begin{equation}\label{eq:sharpness-family}
 \sigma_c=\tfrac14\bigl[\id\otimes\id+c\sigma_z\otimes\id+\sqrt{1-c^2}\,\sigma_x\otimes\sigma_z\bigr],
\end{equation}
an equal mixture of two pure Alice states with Bloch vectors
$(\pm\sqrt{1-c^2},0,c)$ paired with orthogonal Bob states. It is canonical,
separable, $\sigma_c^2=\sigma_c/2$, and $\Dp=p^2/2$. Its correlation matrix
has operator norm $\sqrt{2p-p^2}$, while every local-unitary image of
$\omega_c$ has correlation operator norm $\sqrt p$. Testing against the
norm-one observable $\sigma_x\otimes\sigma_z$ gives
$\norm{\sigma_c-\omega}_1\ge\sqrt{2p-p^2}-\sqrt p=\sqrt p(\sqrt{2-p}-1)$ for
every such $\omega$. Divided by $\Dp^{1/4}=2^{-1/4}\sqrt p$ this tends to
$2^{1/4}(\sqrt2-1)>0$ as $p\downarrow0$, and divided by $\Dp^{\alpha}$ with
$\alpha>1/4$ it diverges.
\end{proof}
The same exponent is optimal for canonical CHSH rigidity: if
$\eps=2(1-c)-M(\sigma)$ then $\Dp\le\eps/2$, so~\eqref{eq:uniform-quarter}
gives an obese state with the same centre vector within
$\min\{2,9(\eps/2)^{1/4}\}$, while~\eqref{eq:sharpness-family} has
$M=1-c^2$ and $\eps=(1-c)^2$.

\section{Full-ellipsoid envelopes}\label{sec:shapes}

The centre length alone does not determine the CHSH optimum of a fixed
ellipsoid. For a fixed canonical state, \eqref{eq:gram} reduces the optimum
over states with that ellipsoid and Bob length $r$ to the choice of the filter
direction; right orthogonal changes of $C$ preserve the range of $C\bv$, so the
optimum depends only on the ellipsoid. We solve this for three classes.

\subsection{Translated spheres}
Let $Q=s^2\id$ and $\cv=ce_z$ with $s,c\ge0$ and $s+c\le1$; the representative
$C=s\diag(1,-1,1)$ is physical.

\begin{theorem}[Translated-sphere envelope]\label{thm:sphere}
For $0\le r<1$ the maximum of $M$ over states with this ellipsoid and Bob
length $r$ is $2s^2+(c^2-s^2)r^2+2scr$; for $s,c,r>0$ it is attained exactly
when $C\bv=sre_z$. For $s>0$ its supremum over $r<1$ is
$2s^2+s^2c^2/(s^2-c^2)$ at $r_*=sc/(s^2-c^2)$ if $sc<s^2-c^2$, and otherwise
$(s+c)^2$, approached as $r\to1$.
\end{theorem}
\begin{proof}
For $sr>0$ let $n=C\bv/(sr)$. By~\eqref{eq:gram},
$TT^{\mathsf T}=q\id+K$ with $q=s^2(1-r^2)$ and
$K=s^2r^2nn^{\mathsf T}+scr(ne_z^{\mathsf T}+e_zn^{\mathsf T})+c^2r^2e_ze_z^{\mathsf T}$,
which is supported on $\operatorname{span}\{n,e_z\}$ with nonpositive
determinant there, so $M=2q+\lambda_{\max}(K)$. For every unit $x$,
$x^{\mathsf T}Kx\le s^2r^2+2scr+c^2r^2$, with equality when $n=x=e_z$, and
for $s,c,r>0$ equality in the cross term forces this alignment. Then
maximize the quadratic in $r$; its vertex lies in $[0,1)$ exactly when
$sc<s^2-c^2$.
\end{proof}
For $s=2/3$, $c=1/3$ the canonical state has $M=8/9<1$, while
$\rho=\frac1{18}\left(\begin{smallmatrix}15&0&0&2\sqrt5\\0&1&0&0\\0&0&0&0\\2\sqrt5&0&0&2\end{smallmatrix}\right)$
has the same ellipsoid and $M=28/27>1$: whitening need not maximize CHSH within
a fixed ellipsoid. The loss-channel boundary $c=1-s$ overlaps Appendix~B
of~\cite{zhang}.

\subsection{Axial ellipsoids}\label{sec:axial}
Take $\cv=ce_z$ and $C=\diag(s,-s,t)$ with $s,t,c\ge0$, and put $S=s^2$,
$T=t^2$, $D=S-T$. The canonical density matrix
$\frac14\left(\begin{smallmatrix}1+c+t&0&0&2s\\0&1+c-t&0&0\\0&0&1-c-t&0\\2s&0&0&1-c+t\end{smallmatrix}\right)$
is positive exactly when
\begin{equation}\label{eq:axial-physical}
 c+t\le1,\qquad 4s^2+c^2\le(1+t)^2.
\end{equation}
The representative $\diag(s,s,t)$ needs $4s^2+c^2\le(1-t)^2$, so whenever it is
physical so is the displayed one.

\begin{theorem}[Axial envelope at fixed Bob length]\label{thm:axial}
For $0\le r<1$,
\begin{equation}\label{eq:axial-envelope}
 \max_{\norm{\bv}=r}M=
 \begin{cases}
 2S+Sr^2(c^2/D-1),& D>0\ \text{and}\ ctr\le D,\\
 S+T+2ctr+(c^2-S)r^2,& D\le0\ \text{or}\ ctr\ge D,
 \end{cases}
\end{equation}
the two expressions agreeing at $ctr=D>0$. For $r>0$ the optimizing filter
directions (any azimuth when a latitude is given) are: $b_z=ctr^2/D$ if $D>0$
and $ctr\le D$; $\bv=re_z$ if $D>0$ and $ctr\ge D$, if $D=0<c$ with $s=t>0$, or
if $D<0<c$; $\bv=\pm re_z$ if $D<0=c$; every direction if $D=c=0$ or $s=t=0$.
\end{theorem}
\begin{proof}
By symmetry $\bv=(\sqrt{R-z^2},0,z)$ with $R=r^2$, $|z|\le r$. The Gram matrix
has the eigenvalue $q=S(1-R)$ and the block
\[
 B(z)=\begin{pmatrix}
 S(1-z^2)&s\sqrt{R-z^2}(c+tz)\\
 s\sqrt{R-z^2}(c+tz)&(1-R)(T-c^2)+(c+tz)^2
 \end{pmatrix},
\]
with $\det(q\id-B)=-S(R-z^2)(1-R)(D+c^2)$. If $D+c^2\ge0$ then
$\lambda_-\le q\le\lambda_+$ and $M=q+\lambda_+$; if $D+c^2<0$ the first
diagonal entry is at least $q$, both block eigenvalues are at least $q$ and
$M=\tr B$ (for $S=0$ the block is diagonal).

\emph{Interior branch} ($D>0$, $ctr\le D$). With $h_i=S+Sc^2R/D$ and
$z_i=ctR/D\in[0,r]$,
\[
 (h_i\id-B)_{11}=Sz^2+Sc^2R/D,\qquad\det(h_i\id-B)=S(D+c^2)(z-z_i)^2.
\]
The first entry is positive unless $z=cr=0$, when the off-diagonal entry
vanishes and the other gap is $S-(1-R)T\ge D>0$. So $h_i\id-B\succeq0$, with
singularity exactly at $z=z_i$.

\emph{Boundary branch.} With $h_b=(t+cr)^2$, $A_b=h_bD+qc^2$ and
$L_b=(c+rt)(t+cr)(ctr-D)$,
\begin{align*}
 (h_b\id-B)_{11}&=-D+2ctr+c^2R+Sz^2,\qquad
 \det(h_b\id-B)=A_b(r-z)^2+2L_b(r-z),\\
 A_br-ct(h_b-q)&=(c+rt)(t+cr)(D-ctr).
\end{align*}
If $D\ge0$ and $ctr\ge D$ the determinant terms are nonnegative, and for
$D>0$ (or $D=0$ with $s=t$, $c,r>0$, where $q>0$) equality holds only at $z=r$.
In the prolate interlacing case $D<0\le D+c^2$ one has $c,t>0$, the first gap
is positive, and with $v=r-z\in[0,2r]$ the determinant is $v(A_bv+2L_b)$,
whose affine factor is $2L_b>0$ at $v=0$ and $2ct(h_b-q)>0$ at $v=2r$; so the
determinant is nonnegative, vanishing only at $z=r$. Finally, if $D+c^2<0$,
$q+(t+cr)^2-M=(r-z)[(T-S)(r+z)+2ct]\ge0$, vanishing only at $z=r$ when $c>0$
and at both poles when $c=0$. The degenerate spheres and points follow from
the block directly.
\end{proof}

\paragraph{Optimizing the Bob length.}
Let $\mathcal M=\sup_{r<1}\max_{\norm{\bv}=r}M$ and $U=S-c^2$. If $D>0$ and
$c^2<D$ then $\mathcal M=2S$, attained only at $r=0$; if $D>0$ and $c^2=D$ the
same value is attained at every radius with $ctr\le D$. Otherwise ($D\le0$ or
$c^2>D$): if $U>ct$ then $\mathcal M=S+ST/U$ at the unique radius
$r_*=ct/U$, with the aligned direction; if $U\le ct$ then
$\mathcal M=S+Sc^2/D$ when $D>0$ and $ct\le D$, and $\mathcal M=(t+c)^2$ when
$D\le0$ or $ct\ge D$, approached as $r\uparrow1$ and not attained except for a
centred line or a point. Indeed the two branches
$F_i(r)=2S+S(c^2-D)r^2/D$ and $F_b(r)=S+T+2ctr-Ur^2$ agree to first order at
$r=D/(ct)$, and when $F_i$ increases the vertex of $F_b$ lies beyond the
transition, since $ct/U-D/(ct)=S(c^2-D)/(ctU)>0$. The boundary value is the
largest squared norm of a point of the ellipsoid, at most one by physicality.

\begin{corollary}[Axial CHSH-breaking criterion]\label{cor:channel}
The physical qubit channel $x\mapsto ce_z+\diag(s,s,t)x$ is CHSH-breaking,
i.e.\ no two-qubit input yields $\beta>2$ after the channel acts on one qubit,
if and only if
\begin{equation}\label{eq:axial-breaking}
 2S\le1\quad\text{and}\quad\bigl[S\le c^2+ct\ \text{or}\ (1-S)(S-c^2)\ge ST\bigr].
\end{equation}
\end{corollary}
\begin{proof}
The canonical state above is the channel's Choi state. Every pure input is a
maximally entangled state filtered on the untouched qubit, and the filter
commutes with the channel, so the outputs of pure inputs with invertible
filters form the orbit optimized above; product inputs (singular filters) give
products, which are CHSH-local, and convexity of $\beta$ reduces mixed inputs
to pure ones. The supremum is $2S$, the vertex value $S+ST/(S-c^2)$ when
$S-c^2>ct$, or a boundary value at most one, and the first two are at most one
exactly under~\eqref{eq:axial-breaking}.
\end{proof}
The criterion concerns arbitrary inputs with one channel use and no further
filtering; filtering on both sides, as in Pal and Ghosh~\cite{pal}, and
Choi-state nonlocality, as in Rutkowski and Siudzi\'nska~\cite{rutkowski}, are
different questions. The unital and amplitude-damping cases are
in~\cite{pal,zhang}.

\subsection{An exact witness requiring a tilted filter}\label{sec:axial-witness}
Take $c=1/5$, $t=3/5$, $S=3/5$, $r=1/2$; \eqref{eq:axial-physical} holds
strictly and $ctr=3/50<D=6/25$. The interior branch gives $M=43/40>1$ at the
canonical-frame filter $\bv=(\sqrt{15}/8,0,1/8)$, while the aligned filter at
the same radius gives $M=47/50<1$. In a rotated Bob basis an attaining state is
\begin{equation}\label{eq:axial-witness}
 \rho_{\mathrm{ax}}=\frac1{160}\begin{pmatrix}
 81&-9\sqrt5&45&15\sqrt5\\
 -9\sqrt5&21&-9\sqrt5&-15\\
 45&-9\sqrt5&39&9\sqrt5\\
 15\sqrt5&-15&9\sqrt5&19
 \end{pmatrix},
\end{equation}
which is positive definite: $\rho_{\mathrm{ax}}=LDL^{\mathsf T}$ with
$D=\diag(81/160,1/10,9/160,1/90)$ and unit lower-triangular
$L$ with rows $(1)$, $(-\sqrt5/9,1)$, $(5/9,-\sqrt5/4,1)$,
$(5\sqrt5/27,-5/12,-\sqrt5/9,1)$. Its Bloch data are $\av=(3/8,0,11/40)$,
$\bv=(0,0,1/2)$ and
$T=\left(\begin{smallmatrix}3\sqrt5/40&0&3/4\\0&-3\sqrt5/10&0\\-9\sqrt5/40&0&1/4\end{smallmatrix}\right)$,
so $\cv=(0,0,1/5)$, $Q=\diag(3/5,3/5,9/25)$ up to rotation and
$T^{\mathsf T}T=\diag(9/32,9/20,5/8)$. The unit settings
$\alpha=(3,0,1)/\sqrt{10}$, $\alpha'=(0,-1,0)$,
$\beta=(0,\sqrt{18},5)/\sqrt{43}$, $\beta'=(0,-\sqrt{18},5)/\sqrt{43}$ give
$\tr(\rho_{\mathrm{ax}}B)=2\sqrt{43/40}$. A second example, $S=1/2$, $t=2/3$,
$c=1/3$, has canonical value $M=1$ and global maximum $M=15/14$ at the aligned
filter with $r=4/7$.

\subsection{Centred ellipsoids}
\begin{theorem}[Centred envelope]\label{thm:centred}
Let $\cv=0$ and let $\ell_1\ge\ell_2\ge\ell_3\ge0$ be the eigenvalues of
$Q=CC^{\mathsf T}$. For $0\le r<1$, $\max_{\norm{\bv}=r}M=\ell_1+(1-r^2)\ell_2$.
For $r>0$ equality holds exactly when $\bv/r$ lies in the top eigenspace of
$C^{\mathsf T}C$ (all directions if $Q=0$).
\end{theorem}
\begin{proof}
With $\bv=rn$, $G=TT^{\mathsf T}=(1-r^2)Q+r^2Cnn^{\mathsf T}C^{\mathsf T}$. For
every rank-two orthogonal projection $P$,
$\tr(PG)\le(1-r^2)(\ell_1+\ell_2)+r^2\norm{Cn}^2\le(1-r^2)(\ell_1+\ell_2)+r^2\ell_1$.
Choosing $n$ in the top right eigenspace and $P$ containing $Cn$ and a second
leading left eigenvector attains both bounds; for $0<r<1$ equality forces
$\norm{Cn}^2=\ell_1$.
\end{proof}
Over all radii the supremum is the canonical value $\ell_1+\ell_2$. The
corresponding CHSH-breaking criterion for unital qubit channels is prior
work~\cite{pal,zhang}; the fixed-radius value and equality cases are the
finer statement here.

\section{A printed optimizer in Zhang et al.}\label{sec:zhang}
Appendix~B of~\cite{zhang} (arXiv~v2 and the published version) prints an
optimizing vector with a repeated $R_{31}$ and last coordinate
$R_{33}-(1-\alpha^2)$, and its prose states $R_{33}=1-\alpha^2$. Re-expanding
its preceding correlation matrix gives, for $r=\sqrt{1-\alpha^2}$,
$H'=-p(1-p)aa^{\mathsf T}+r^2e_ze_z^{\mathsf T}$ with
$a=(\alpha R_{31},\alpha R_{32},R_{33}-r)$, so the optimizing condition is
$R_{33}=r$, which is $b_z=\norm{\bv}^2$ in filter coordinates. For $p=1/2$,
$r=1/2$, $\alpha=\sqrt3/2$ and a $y$ rotation with $R_{33}=1/4$, the correlation
matrix is
$T=\left(\begin{smallmatrix}\sqrt6/16&0&-3\sqrt5/16\\0&-\sqrt6/4&0\\\sqrt{30}/8&0&3/8\end{smallmatrix}\right)$
and $M=(399+9\sqrt{145})/512<1$, strictly, since $113^2-81\cdot145=1024>0$;
the printed direction does not attain the maximum, which the corrected
direction does. The maximal value and the CHSH-breaking threshold
of~\cite{zhang} are correct.

\section{Verification}\label{sec:formal}

\subsection{Lean}
The package contains a standalone Lean project (Lean~4.19.0, Mathlib
revision \texttt{c44e0c8ee63ca166\allowbreak450922a373c7409c\allowbreak5d26b00b}) of fifteen modules and a
whole-project axiom audit. Its definitions are explicit: complex $4\times4$
matrices; density matrices as Hermitian, trace one and
$\operatorname{Re}\langle w,\rho w\rangle\ge0$ for \emph{every} complex vector
$w$; the standard Pauli matrices and the tensor product in the basis
$\ket{00},\ket{01},\ket{10},\ket{11}$; the Bell operator above with four real
unit vectors; the modulus $\sqrt{\operatorname{normSq}}$ of the complex
expectation $\tr(\rho B)$; and $\beta$ as the real supremum of these values.
\begin{itemize}
\item \texttt{CertifiedCHSH.physical\_joint\_chsh\_attainment}: for every
$0\le c\le1$ and $r^2<1$ the matrix $\rho^\sharp_{c,r}$ of
Proposition~\ref{prop:dyadic} is a density matrix, the centre formula
\eqref{eq:centre} applied to its Bloch data gives $(0,0,c)$, Bob's length is
$r$, and $\beta(\rho^\sharp_{c,r})=2\sqrt{1+(1-2c)(1-r^2)}$. The Bloch data are
tied to the matrix by the readout theorems
\texttt{CHSHPhysical.alice\_readout}, \texttt{bob\_readout} and
\texttt{correlation\_readout}. The upper bound covers every setting through an
elementary orthogonal-plane certificate, and explicit unit settings attain it;
the Horodecki formula is not assumed.
\item \texttt{CHSHAxialWitness.beta\_exact} and related theorems:
the matrix~\eqref{eq:axial-witness} is a density matrix with the stated Pauli
data and $\beta=2\sqrt{43/40}$.
\item \texttt{BoundarySharpness}: the family~\eqref{eq:sharpness-family} is a
canonical density matrix with $\rho^2=\rho/2$, purity $1/2$, the stated Pauli
data, exact $\beta$ and canonical CHSH gap $(1-c)^2$.
\item \texttt{AxialCertificates.axial\_exact\_maximum}: for $S,c,t,r\ge0$,
$r^2\le1$ and $|z|\le r$, the largest sum of two of the three roots of the
axial Gram matrix (the explicit radical roots of the block and $q$) is at most
the envelope~\eqref{eq:axial-envelope}, with an attaining $z$; neither
physicality nor the inequality is assumed. Further modules check the obese
block algebra, the orientation correction, spectral purity inequalities and
the countercertificate of Section~\ref{sec:zhang}.
\end{itemize}
A fresh replay compiled all modules with warnings as errors and audited all
839 declarations owned by the fifteen modules: the 543 declarations other than
code-generation auxiliaries are not axioms and depend only on
\texttt{propext}, \texttt{Classical.choice} and \texttt{Quot.sound}; the 296
auxiliaries stored by the Lean~4.19 code generator are recorded separately,
and one of them refers to the generator's placeholder \texttt{lcProof}. The sources contain no
\texttt{sorry}, added axiom or \texttt{native\_decide}.

\emph{Not formalized:} the arbitrary-state bound~\eqref{eq:H} (it is prior
work); the general filter-to-Gram correspondence~\eqref{eq:gram} and the
Horodecki identification for arbitrary states, which connect the theorem
\texttt{axial\_\allowbreak exact\_\allowbreak maximum} to Theorem~\ref{thm:axial}; the equality
classification; the rigidity theorems and the lower bound against arbitrary
local-unitary images; and the operational meaning of the centre.

\subsection{Exact and numerical checks}
An exact SymPy program checks 46 identities: the Bloch data, centre, Gram
characteristic polynomial and $M=H$ for $\rho^\sharp_{c,r}$ symbolically; the
block identities~\eqref{eq:factor} and their agreement with~\eqref{eq:gram};
the Gram formula on an explicitly filtered obese state; all axial identities
in the proof of Theorem~\ref{thm:axial}; the axial witness (positivity,
readouts, whitening, shape and Bell value); the sharpness family; the
countercertificate of Section~\ref{sec:zhang}; and the examples $\rho_*$ and
the translated sphere. A seeded NumPy program searches for counterexamples:
$10^4$ random states of ranks one to four against~\eqref{eq:H}
and~\eqref{eq:shape-gap}; the equality orbit and
Proposition~\ref{prop:orientation}; brute-force filter directions for 300
random axial shapes against~\eqref{eq:axial-envelope}; random centred
ellipsoids against Theorem~\ref{thm:centred}; the channel criterion against a
radius grid; and the sharpness lower bound. No violation beyond rounding
($\sim10^{-15}$) was found. These searches are falsification attempts, not
proofs.

\section{Limitations}\label{sec:limits}
The exact optimum for an unrestricted displaced triaxial ellipsoid remains
open; \eqref{eq:gram} reduces it to an optimization over Bob's Bloch ball for
each measurement plane, which is not solved here. The shape-gap
bound~\eqref{eq:shape-gap} is not sharp for general ellipsoids. The rigidity
constants are not claimed sharp; the uniform exponent is. The full text of
Xing and Yang~\cite{xing}, on volume-based bounds, was not available, so that
paper is not excluded as an overlapping source.

\section{Research process and AI disclosure}
This is fully autonomous agentic research. Autonomous AI research agents
produced the mathematics, the Lean development, the computations, an audit of
the public Omega proof and an earlier version of this text in earlier
sessions; their records do not name the system used, and the reviews in those
records were AI-assisted reviews within the same workflow. Opus 5.5 by
Anthropic used with Claude Code assessed the package for this release: it
compared the results with Conjecture~1, the Omega dossier and source, and the
2020 paper, and determined what is new; read the central proofs; wrote and ran
its own exact and numerical check programs; replayed the Lean project with a
whole-project axiom audit; rewrote this manuscript; and integrated the package
into the repository. The audit and replay records of the third-party Omega
project are not included. Dev D. Goyal is responsible for publication. The
work has not been externally peer reviewed. The formal scope and remaining
limits are stated in Sections~\ref{sec:formal} and~\ref{sec:limits}.

\begingroup\small
\begin{thebibliography}{10}
\bibitem{milne}
A.~Milne, D.~Jennings, S.~Jevtic and T.~Rudolph,
Quantum correlations of two-qubit states with one maximally mixed marginal,
Phys.\ Rev.\ A \textbf{90}, 024302 (2014),
\href{https://arxiv.org/abs/1404.3951}{arXiv:1404.3951}.
\bibitem{dossier}
The Omega Institute, CHSH nonlocality is tightly bounded by the
steering-ellipsoid centre, problem dossier and Lean source, repository
\texttt{the-omega-institute/trureturing}, issue \#16558; dossier committed
14 October 2026 (commit \texttt{af41108}),
\url{https://github.com/the-omega-institute/trureturing/blob/af41108edc85c119837189f120c6d2d9aa3f5a25/Problems/milne-jennings-jevtic-rudolph-2014-chsh-centre-bound.md};
source read at commit \texttt{609e95e},
\url{https://github.com/the-omega-institute/trureturing/blob/609e95e447fb1045172423c4d7bcf2a6e321eef0/D5/S3/Quantum/Entanglement/SteeringEllipsoidCHSH.lean}.
Apache-2.0.
\bibitem{zhang}
Y.~Zhang, R.~Araiza Bravo, V.~O.~Lorenz and E.~Chitambar,
Channel activation of CHSH nonlocality, New J.\ Phys.\ \textbf{22}, 043003
(2020), Appendix~B,
\href{https://doi.org/10.1088/1367-2630/ab7bef}{doi:10.1088/1367-2630/ab7bef},
\href{https://arxiv.org/abs/1911.06349}{arXiv:1911.06349}.
\bibitem{ducuara}
A.~F.~Ducuara, C.~E.~Susa and P.~Skrzypczyk,
Locally inaccessible hidden quantum correlations, Phys.\ Rev.\ A \textbf{110},
022435 (2024), \href{https://arxiv.org/abs/2406.03945}{arXiv:2406.03945}.
\bibitem{horodecki}
R.~Horodecki, P.~Horodecki and M.~Horodecki,
Violating Bell inequality by mixed spin-$\frac12$ states: necessary and
sufficient condition, Phys.\ Lett.\ A \textbf{200}, 340--344 (1995),
\href{https://doi.org/10.1016/0375-9601(95)00214-N}{doi:10.1016/0375-9601(95)00214-N}.
\bibitem{pal}
R.~Pal and S.~Ghosh, Non-locality breaking qubit channels: the case for CHSH
inequality, J.\ Phys.\ A \textbf{48}, 155302 (2015),
\href{https://arxiv.org/abs/1306.3151}{arXiv:1306.3151}.
\bibitem{rutkowski}
A.~Rutkowski and K.~Siudzi\'nska, Controlling nonlocality of bipartite qubit
states via quantum channels, Phys.\ Lett.\ A (2025) 130313,
\href{https://doi.org/10.1016/j.physleta.2025.130313}{doi:10.1016/j.physleta.2025.130313},
\href{https://arxiv.org/abs/2407.16035}{arXiv:2407.16035}.
\bibitem{xing}
L.-L.~Xing and H.~Yang, Bounding entanglement and Bell nonlocality for
arbitrary two-qubit state via quantum steering ellipsoid, Phys.\ Scr.\
\textbf{101}, 135103 (2026),
\href{https://doi.org/10.1088/1402-4896/ae5354}{doi:10.1088/1402-4896/ae5354}.
\end{thebibliography}
\endgroup
\end{document}
